\chapter{Supplementary Optimization Background}

\section{Oracle Tiers: First- and Second-Order Methods}
\label{app:oracle-tiers}

The three tiers of oracle access form an information hierarchy that
governs both the choice and the efficiency of applicable algorithms.

\paragraph{First-Order Methods.}
First-order methods additionally access the gradient
$\nabla f(x) \in \mathbb{R}^d$, which encodes the direction of
steepest ascent at $x$. For a differentiable function
$f : \mathbb{R}^d \to \mathbb{R}$, gradient-based methods exploit
this to define descent directions and form the backbone of classical
continuous optimization~\cite{nocedal2006numerical}. When $f$ is
convex but non-differentiable, the gradient generalizes to the
\textit{subdifferential} $\partial f(x)$, the set of all
\textit{subgradients} $g$ satisfying
\[
  f(y) \geq f(x) + g^\top(y - x) \quad \forall\, y \in \mathbb{R}^d.
\]
Subgradient methods replace the gradient with any element of
$\partial f(x)$, recovering convergence under convexity alone, albeit
at rates slower than those attainable with smooth
gradients~\cite{bertsekas2009convex}.

\paragraph{Second-Order Methods.}
Second-order methods further access the Hessian
$\nabla^2 f(x) \in \mathbb{R}^{d \times d}$, which captures local
curvature. Methods that utilize or approximate the Hessian, such as
Newton and quasi-Newton schemes, can achieve super-linear or quadratic
convergence near a solution~\cite{nocedal2006numerical}, at the cost
of maintaining a $d \times d$ matrix. Together, the three tiers
define distinct levels of oracle access; restricting access
progressively from second-order to zeroth-order culminates in the
black-box model where only function values are observable.

\section{Landscape Properties: Convexity and Smoothness}
\label{app:landscape-props}

\paragraph{Convexity.}
A function $f : \mathbb{R}^d \to \mathbb{R}$ is \textit{convex} if, for all
$x,y \in \mathbb{R}^d$ and $\lambda \in [0,1]$,
\[
  f\!\left(\lambda x + (1-\lambda)y\right)
  \leq \lambda f(x) + (1-\lambda)f(y)
\]
\cite[Def.~1.1.1]{bertsekas2009convex}. For a convex function over a convex
domain, every local minimum is also a global minimum
\cite[Prop.~1.3.2]{bertsekas2009convex}. A differentiable function is
\textit{strongly convex} with parameter $\mu>0$ if
\[
  f(y) \geq f(x) + \nabla f(x)^\top (y-x)
  + \tfrac{\mu}{2}\|y-x\|^2,
\]
which is one of the standard equivalent characterizations of strong convexity
\cite[Thm.~2.1.9]{nesterov2018lectures}. Strong convexity implies a unique
global minimizer \cite[Thm.~2.1.8]{nesterov2018lectures}. Moreover, if $f$
is also $L$-smooth, then gradient descent with a suitable fixed step size
converges linearly \cite[Thm.~2.1.15]{nesterov2018lectures}.

\paragraph{Smoothness and Lipschitz Continuity.}
A differentiable function $f$ is \textit{$L$-smooth} if its gradient
is Lipschitz continuous with constant $L > 0$:
\[
  \|\nabla f(x) - \nabla f(y)\| \leq L\|x - y\|
  \quad \forall\, x, y \in \mathbb{R}^d.
\]
Smoothness bounds the curvature of $f$ from above and is the standard
assumption under which gradient descent achieves an $O(1/T)$
convergence rate for convex objectives~\cite{nocedal2006numerical}.
The objective $f$ itself is \textit{Lipschitz continuous} with
constant $G$ if $|f(x) - f(y)| \leq G\|x - y\|$, a weaker condition
that enables subgradient convergence without differentiability.
Neither property is assumed in the black-box pseudo-Boolean setting
studied in this project.

\section{Haj\'{e}k's Convergence Condition for Simulated Annealing}
\label{app:hajek}

Modelling simulated annealing as a time-inhomogeneous Markov chain
on the state space $\mathcal{X}$, Haj\'{e}k~\cite{hajek1988cooling}
established the following sharp necessary and sufficient condition for
convergence in probability to the global optimum.

\begin{theorem}[Haj\'{e}k 1988]
  Let $d^{*}$ denote the \emph{depth} of the deepest local minimum
  that is not globally optimal, defined as the minimum energy that
  must be crossed to escape it. A cooling schedule $T(t) > 0$
  guarantees that the distribution of the current state converges to
  a distribution supported on the global optima if and only if
  \[
    \sum_{t=1}^{\infty}
    \exp\!\left(\frac{-d^{*}}{T(t)}\right) = +\infty.
  \]
\end{theorem}

For the logarithmic schedule $T(t) = d / \log(t + 1)$, the condition
reduces to the requirement $d \geq d^{*}$: the scaling constant must
be at least as large as the hardest energy barrier in the
landscape~\cite{hajek1988cooling,bertsimas1993simulated}. Logarithmic
cooling is therefore theoretically optimal in the sense that it is
the fastest schedule satisfying the convergence condition.

In practice, $d^{*}$ is unknown for general combinatorial problems,
and logarithmic cooling is often too slow to be useful within a
finite evaluation budget. Exponential cooling
$T(t) = T_0 \cdot \alpha^t$ ($\alpha \in (0,1)$) and linear cooling
$T(t) = \max(T_f,\, T_0 - (T_0 - T_f)\,t/t_{\max})$ cool faster but
violate the condition once the temperature reaches zero, at which
point \acrshort{sa} reduces to a deterministic hill-climber. The
choice of schedule therefore encodes an explicit trade-off between
the convergence guarantee and the practical computational budget.

\section{Plateau Function: Runtime Analysis}
\label{app:plateau-runtime}

\citet{antipov2021precise} establish that for any unbiased mutation
operator the expected runtime of the $(1+1)$-\acrshort{ea} on
\textsc{Plateau}$_k$ is
\[
  \mathbb{E}[T]
  = \frac{n^k}{k!\,\Pr[1 \leq \alpha \leq k]}\,(1 + o(1)),
\]
where $\alpha$ is the number of bits flipped per step. This
expression depends only on the probability of flipping between one
and $k$ bits, not on any other property of the operator. The optimal
mutation rate under standard bit mutation is approximately $k/(en)$,
a factor of $\Theta(k)$ above the canonical $1/n$ recommendation,
mirroring the intuition that wider plateaus require more aggressive
perturbation to exit. This result is directly relevant to the suite
06 and suite 10 experimental analyses, where the contrast between
Jump and Plateau regret trajectories is examined at shared $k$ values.
